\section{Related work}
\label{sec:related}

\paragraph{Differences between documents.}
Line-based differences go back to Hunt and McIlroy~\cite{hunt1976} and Myers's
$O(ND)$ algorithm~\cite{myers1986}, which Git uses by default.
\texttt{latexdiff}~\cite{latexdiff} marks up the differences between two
\LaTeX{} sources so that they can be typeset together. Papilot instead
typesets each version separately and shows the corresponding strips of the
two PDFs one above the other, which keeps the typeset result of each version
exact at the cost of a second build.

\paragraph{Concurrent editing.}
Operational transformation~\cite{ellis1989} and conflict-free replicated data
types~\cite{shapiro2011} let several writers edit a shared document and converge
without coordination. Papilot's setting is different: proposals are not applied
until an author accepts them, so convergence is a property of accepted edits
only, and \cref{lem:tp1} needs to hold only for compatible pairs. Incompatible
pairs are surfaced to the author rather than transformed.

\paragraph{Line breaking and text layers.}
The Knuth--Plass algorithm~\cite{knuth1981} is the reason \LaTeX{} spaces vary
from line to line, and with it the drift analysed in \cref{sec:geometry}. PDF
viewers such as PDF.js~\cite{pdfjs} place an invisible text layer over the
rendered page; Papilot uses PDF.js for rendering, search and links, and
replaces only the placement of its text with the per-word runs of
\cref{thm:per-word}.

\paragraph{Source synchronisation.}
SyncTeX~\cite{laurens2008} maps between typeset positions and source lines and
underlies forward and inverse search in most \TeX{} editors. Papilot uses its
reverse queries and corrects their line-boundary error with
\cref{alg:settle}.

\paragraph{Queueing.}
The waiting-time results of \cref{sec:agents} are classical; we use Erlang's
formula~\cite{erlang1917}, Little's law~\cite{little1961} and the treatment of
$M/M/c$ queues in Kleinrock~\cite{kleinrock1975}.
